\begin{examplebox}
\textbf{例 1：由定义写方程。} $F_1(-3,0),F_2(3,0)$，$PF_1+PF_2=10$。
先查 $10>6$，再读出 $a=5,c=3$，故 $b^2=25-9=16$。
轨迹为 $\boxed{x^2/25+y^2/16=1}$。
回验：短轴端点 $(0,4)$ 到两焦点的距离各为 $5$，和确为 $10$。

\medskip
\textbf{例 2：求一条焦半径，先用减法。}
在 $x^2/25+y^2/9=1$ 上，若 $PF_1=4$，则 $a=5,c=4$，
所以 $PF_2=2a-4=\boxed{6}$，无须先求点坐标。
回验：$r_1-r_2=2cx/a$ 给出 $x=-5/4$；代入椭圆得 $y^2=135/16$，
确有 $PF_1^2=16,PF_2^2=36$，满足条件的点实际存在。
\end{examplebox}
\clearpage
\begin{examplebox}
\textbf{例 3：焦点三角形的面积。}
在 $x^2/25+y^2/16=1$ 上，$\angle F_1PF_2=60^\circ$，求 $\triangle PF_1F_2$ 的面积。
\begin{center}\begin{tikzpicture}[font=\small,>=stealth]
\draw[black!30,rounded corners=2pt] (0,0) rectangle (14.4,8.2);
\node[anchor=west,font=\small\bfseries] at (.5,7.72) {例 3：距离和与余弦定理联立，不必分别求两边};
\begin{scope}[shift={(7.2,3.75)},scale=.58]
\draw[blue,very thick] (0,0) ellipse (5 and 4);
\draw[black!55] (-3,0)--(3,0);
\draw[orange!85!black,very thick] (-3,0)--(3.1914236925,3.0792014357) node[midway,above left] {$m$};
\draw[violet,very thick] (3,0)--(3.1914236925,3.0792014357) node[midway,left] {$n$};
\fill (-3,0) circle (2pt) node[below] {$F_1$};
\fill (3,0) circle (2pt) node[below] {$F_2$};
\fill (3.1914236925,3.0792014357) circle (2pt) node[above right] {$P$};
\node[below] at (0,0) {$6$};
\end{scope}
\node at (7.2,6.7) {$\angle F_1PF_2=60^\circ,\quad m+n=10$};

\node[anchor=west] at (.7,.55) {蓝色椭圆 $x^2/25+y^2/16=1$；三角形面积也可用底 $6$ 与高 $y_P$ 回验};
\end{tikzpicture}\end{center}

$a=5,c=3$。记 $m=PF_1,n=PF_2$，由定义 $m+n=10$。
余弦定理给出 $36=m^2+n^2-mn=100-3mn$，所以 $mn=64/3$。
\[
S=\frac12mn\sin60^\circ=\boxed{\frac{16\sqrt3}{3}}.
\]
\textbf{坐标回验。} $m,n=5\pm\sqrt{33}/3$，可取
$P(5\sqrt{33}/9,16\sqrt3/9)$，代入椭圆为 $1$，
用底乘高得 $S=\tfrac12\cdot6\cdot16\sqrt3/9=16\sqrt3/3$，与前法一致。
\end{examplebox}
\clearpage
